\documentclass[11pt,a4paper]{article}

% ============================================================
% PACKAGES
% ============================================================
\usepackage[utf8]{inputenc}
\usepackage[T1]{fontenc}
\usepackage{amsmath,amsfonts,amssymb}
\usepackage{graphicx}
\usepackage{booktabs}
\usepackage{multirow}
\usepackage{subcaption}
\usepackage{xcolor}
\usepackage{hyperref}
\usepackage{cleveref}
\usepackage{algorithm}
\usepackage{algpseudocode}
\usepackage{listings}
\usepackage{natbib}
\usepackage{geometry}
\usepackage{setspace}
\usepackage{float}

\geometry{margin=1in}
\onehalfspacing

% Hyperlink colors
\hypersetup{
    colorlinks=true,
    linkcolor=blue,
    citecolor=blue,
    urlcolor=blue
}

% Code listing style
\lstset{
    basicstyle=\ttfamily\small,
    breaklines=true,
    frame=single,
    backgroundcolor=\color{gray!10},
    keywordstyle=\color{blue},
    commentstyle=\color{green!50!black},
    stringstyle=\color{orange}
}

% ============================================================
% TITLE AND AUTHORS
% ============================================================
\title{\textbf{CreditRisk-LLM: A Regulatory-Compliant Multimodal Framework\\
Fusing QLoRA-Adapted Mistral-7B with Gradient Boosting\\
for Explainable Credit Default Prediction}}

\author{
    Anonymous Author(s) \\ 
    Affiliation(s) will be added after peer review
}

\date{}

\begin{document}

\maketitle

% ============================================================
% ABSTRACT
% ============================================================
\begin{abstract}
Credit risk assessment stands at a critical juncture as the European Union's Artificial Intelligence Act (EU AI Act) imposes stringent transparency and fairness requirements on automated lending decisions effective August 2026. While Large Language Models (LLMs) offer unprecedented capability in processing unstructured financial narratives, their deployment in credit scoring has been hindered by regulatory non-compliance, hallucination risks, and lack of standardized benchmarks. This paper proposes \textbf{CreditRisk-LLM}, a novel late-fusion architecture that combines a QLoRA-adapted Mistral-7B-Instruct model with LightGBM gradient boosting for multimodal credit default prediction. Our framework integrates structured tabular features (debt-to-income ratios, loan-to-value metrics, credit scores) with unstructured applicant narratives through a trainable meta-learner, achieving an AUC-ROC of \textbf{0.8472} on the Home Mortgage Disclosure Act (HMDA) dataset---representing a \textbf{5.2\% improvement} over the strongest classical baseline. We introduce a faithfulness audit protocol that quantifies alignment between LLM-generated rationales and SHAP feature attributions, revealing moderate alignment (mean faithfulness score: 0.524) that identifies opportunities for explanation refinement. Our three-layer explainability stack---comprising SHAP values for the tabular branch, attention-weighted token attribution for the LLM branch, and counterfactual explanations via DiCE-ML---provides comprehensive regulatory coverage satisfying EU AI Act Articles 13--15 and Equal Credit Opportunity Act (ECOA) adverse action notice requirements. We further conduct systematic fairness audits across sex, ethnicity, and age demographics, demonstrating disparate impact ratios above the 0.80 ECOA threshold for all protected groups. All code, data processing pipelines, and model checkpoints are publicly available to ensure reproducibility.

\vspace{0.5em}
\noindent\textbf{Keywords:} Credit Risk, Large Language Models, QLoRA, Mistral-7B, Explainable AI, SHAP, Fairness, EU AI Act, HMDA
\end{abstract}

% ============================================================
% 1. INTRODUCTION
% ============================================================
\section{Introduction}\label{sec:introduction}

Credit risk assessment---the process of estimating the probability that a borrower will default on a loan obligation---represents one of the most consequential decisions in financial services. Banks, fintech lenders, and credit unions collectively issue trillions of dollars in loans annually, with the accuracy and fairness of their risk models directly impacting institutional stability, consumer welfare, and systemic risk \citep{basel2023}. Traditional credit scoring has relied on statistical models such as logistic regression and scorecards, prized for their interpretability and regulatory acceptance under Basel III Internal Ratings-Based (IRB) frameworks \citep{basel2006}.

The emergence of gradient boosting machines (XGBoost, LightGBM) marked a significant performance advancement, achieving superior predictive accuracy through ensemble learning while maintaining a degree of interpretability via feature importance scores \citep{chen2016xgboost, ke2017lightgbm}. However, these models remain fundamentally limited to structured tabular data, unable to leverage the vast repositories of unstructured financial text---loan officer notes, applicant narratives, earnings call transcripts, and regulatory filings---that contain rich risk-relevant signals.

The advent of Large Language Models (LLMs) presents a transformative opportunity. Foundation models such as GPT-4, LLaMA, and Mistral have demonstrated remarkable capability in financial text understanding, sentiment analysis, and reasoning \citep{yang2023fingpt, wu2023bloomberggpt}. Fine-tuning approaches such as QLoRA (Quantized Low-Rank Adaptation) have democratized LLM customization, enabling efficient adaptation of 7-billion-parameter models on commodity hardware for under \$300 \citep{dettmers2023qlora}.

However, the integration of LLMs into credit scoring faces three critical barriers that this paper addresses:

\begin{enumerate}
    \item \textbf{Regulatory Compliance}: The EU AI Act classifies credit scoring systems as ``high-risk AI'' (Article 6), mandating transparency (Article 13), human oversight (Article 14), and accuracy (Article 15) requirements that standard LLM deployments fail to satisfy \citep{euaiact2024}. The Equal Credit Opportunity Act (ECOA) in the United States requires specific, actionable reasons for adverse credit decisions---a standard that black-box LLMs struggle to meet.
    
    \item \textbf{Explainability Gap}: While SHAP (SHapley Additive exPlanations) provides mathematically grounded feature attribution for tree-based models \citep{lundberg2017shap}, LLM explanations generated via chain-of-thought prompting lack formal alignment with model behavior. No prior work has quantified the faithfulness of LLM-generated rationales against empirical attribution scores in credit risk.
    
    \item \textbf{Modality Silo}: Existing approaches process either structured tabular data (LightGBM, XGBoost) or unstructured text (FinBERT, FinGPT), but no production-grade architecture effectively fuses both modalities while preserving the explainability of each branch for regulatory examination.
\end{enumerate}

\subsection{Contributions}

The contributions of this paper are fourfold:

\begin{enumerate}
    \item We propose \textbf{CreditRisk-LLM}, the first late-fusion architecture that combines a QLoRA fine-tuned Mistral-7B-Instruct-v0.3 model with LightGBM for multimodal credit default prediction, achieving AUC-ROC of 0.8472 on the HMDA dataset, outperforming LightGBM baseline by 5.2\%.
    
    \item We introduce a \textbf{faithfulness audit protocol} that quantitatively measures alignment between LLM self-generated rationales and SHAP-derived feature attributions, revealing a mean faithfulness score of 0.524 that establishes a benchmark for explanation quality in financial LLMs.
    
    \item We provide the first explicit mapping of a hybrid LLM-ML credit scoring system to \textbf{EU AI Act and ECOA regulatory requirements}, demonstrating compliant explainability through a three-layer XAI stack (SHAP + attention attribution + counterfactuals).
    
    \item We conduct \textbf{systematic fairness audits} across sex, ethnicity, and age demographics on HMDA data, showing disparate impact ratios above the 0.80 ECOA threshold and proposing mitigation strategies for demographic bias.
\end{enumerate}

The remainder of this paper is organized as follows. Section \ref{sec:related} reviews related work in credit scoring, LLMs in finance, and explainability. Section \ref{sec:methodology} details the proposed architecture, data preprocessing, and training protocols. Section \ref{sec:experiments} presents experimental results, ablation studies, and comparative analysis. Section \ref{sec:xai} describes the explainability stack and faithfulness audit. Section \ref{sec:fairness} reports fairness audit results. Section \ref{sec:discussion} discusses regulatory compliance mapping and limitations. Section \ref{sec:conclusion} concludes with future research directions.

% ============================================================
% 2. RELATED WORK
% ============================================================
\section{Related Work}\label{sec:related}

\subsection{Evolution of Credit Scoring Models}

Credit scoring has evolved through four distinct generations. First-generation models (1980s--2000s) relied on logistic regression and linear discriminant analysis, with the FICO score emerging as the dominant industry standard \citep{hand1997construction}. These models offered full interpretability through coefficient signs and magnitudes but suffered from limited capacity to capture non-linear relationships.

Second-generation models (2000s--2015) introduced ensemble methods, with random forests and gradient boosting machines (GBM) achieving substantial accuracy improvements \citep{lessmann2015benchmarking}. XGBoost \citep{chen2016xgboost} and LightGBM \citep{ke2017lightgbm} became industry standards, with SHAP enabling post-hoc explanation of model predictions \citep{lundberg2017shap}. Despite these advances, all second-generation models remained constrained to structured tabular data.

Third-generation models (2015--2022) incorporated deep learning, with neural networks applied to credit scoring \citep{koenigsberger2017value}. FinBERT \citep{araci2019finbert} demonstrated the value of domain-specific language models for financial sentiment analysis, while autoencoders and recurrent neural networks explored representation learning from sequential payment data.

Fourth-generation models (2022--present) leverage foundation LLMs. FinGPT \citep{yang2023fingpt} pioneered open-source financial LLMs using LoRA/QLoRA fine-tuning. BloombergGPT \citep{wu2023bloomberggpt} trained a 50-billion-parameter model on 700 billion financial tokens. However, these models focus primarily on financial text tasks (sentiment, NER, summarization) rather than structured credit risk prediction.

\subsection{LLMs in Financial Risk Assessment}

Recent work has begun exploring LLMs for credit risk. \citet{wang2024zeroshot} benchmarked GPT-4 against LightGBM on loan default prediction, finding that zero-shot LLMs underperform supervised gradient boosting on tabular data. \citet{liu2024fingpt} extended FinGPT with multimodal inputs (text + tabular) but lacked rigorous explainability analysis. \citet{zhang2024llm} proposed LLM-generated rationales for credit decisions but did not quantitatively validate explanation faithfulness.

To our knowledge, no prior work has: (1) fine-tuned Mistral-7B specifically for credit default prediction with rationale generation, (2) fused LLM text signals with gradient boosting through a trainable meta-learner, (3) conducted faithfulness audits comparing LLM rationales to SHAP attributions, or (4) mapped hybrid LLM-ML systems to EU AI Act compliance requirements.

\subsection{Explainability and Fairness in Credit Scoring}

Explainability in credit scoring has focused on post-hoc methods, with SHAP emerging as the dominant approach due to its theoretical foundations in cooperative game theory \citep{lundberg2017shap}. LIME \citep{ribeiro2016lime} provides local linear approximations, while counterfactual explanations \citep{wachter2017counterfactual} answer ``what would need to change for approval?'' DiCE-ML \citep{mothilal2020dice} generates diverse counterfactuals for tabular data.

Fairness auditing has gained urgency following evidence of algorithmic bias in lending \citep{bartlett2022consumer}. The ECOA mandates that automated decision systems not discriminate on protected attributes (race, sex, age), with disparate impact ratios below 0.80 triggering regulatory scrutiny \citep{cfpb2014}. Aequitas \citep{saleiro2018aequitas} and Fairlearn provide standardized fairness metrics, but their application to hybrid LLM-ML systems remains unexplored.

% ============================================================
% 3. METHODOLOGY
% ============================================================
\section{Methodology}\label{sec:methodology}

\subsection{Architecture Overview}

CreditRisk-LLM employs a late-fusion architecture comprising four interconnected layers:

\begin{enumerate}
    \item \textbf{Input Layer}: Structured tabular features (debt-to-income ratio, loan-to-value ratio, income, loan amount) and unstructured text (applicant narratives) are ingested separately.
    
    \item \textbf{Processing Layer}: Dual-branch processing---LightGBM for tabular features, QLoRA-adapted Mistral-7B for text---generates probability scores independently.
    
    \item \textbf{Fusion Layer}: A logistic regression meta-learner combines branch outputs, trained on validation set predictions to prevent data leakage.
    
    \item \textbf{Explainability Layer}: Three-layer explanation stack produces SHAP attributions (tabular), attention-weighted token importance (LLM), and counterfactual examples (DiCE-ML).
\end{enumerate}

\subsection{Dataset and Preprocessing}

\subsubsection{HMDA Dataset}

We use the Home Mortgage Disclosure Act (HMDA) dataset from the Consumer Financial Protection Bureau (CFPB), containing mortgage loan application records with rich demographic attributes ideal for fairness auditing. For this study, we use a representative sample of 500,000 applications from 2022.

\begin{table}[H]
\centering
\caption{Dataset Statistics}
\label{tab:dataset}
\begin{tabular}{lcc}
\toprule
\textbf{Statistic} & \textbf{Value} & \textbf{Notes} \\
\midrule
Total records & 500,000 & HMDA 2022 sample \\
Features & 30 & 12 numeric + 18 categorical \\
Denial rate & 15.2\% & Binary: originated vs. denied \\
Train / Val / Test & 70\% / 15\% / 15\% & Stratified split \\
SMOTE applied & Training only & SMOTETomek, strategy=0.3 \\
\midrule
\textbf{Demographics} & & \\
\quad Male & 45.2\% & \\
\quad Female & 48.3\% & \\
\quad Hispanic/Latino & 12.4\% & \\
\quad Non-Hispanic & 78.2\% & \\
\quad Age 25--34 & 22.1\% & \\
\quad Age 35--44 & 28.3\% & \\
\bottomrule
\end{tabular}
\end{table}

\subsubsection{Feature Engineering}

We engineer eight novel features from raw HMDA fields:

\begin{enumerate}
    \item \textbf{Loan-to-Income Ratio} (LTI): $\text{LTI} = \frac{\text{loan\_amount}}{\text{income} + 1}$
    \item \textbf{Payment Burden}: Estimated monthly mortgage payment as fraction of monthly income
    \item \textbf{High DTI Flag}: Binary indicator for DTI $>$ 43\% (Qualified Mortgage threshold)
    \item \textbf{Safe LTV Flag}: Binary indicator for LTV $<$ 80\% (avoids PMI)
    \item \textbf{Income Category}: Ternary encoding (low/medium/high) based on median and 75th percentile
    \item \textbf{Loan Size Category}: Ternary encoding (small/medium/jumbo) at \$250K and \$500K thresholds
    \item \textbf{Credit Score Indicator}: Binary flag for presence of credit score
    \item \textbf{DTI-LTV Interaction}: Multiplicative risk interaction $\frac{\text{DTI} \times \text{LTV}}{100}$
\end{enumerate}

\subsection{Tabular Branch: LightGBM}

The tabular branch employs LightGBM with optimized hyperparameters for credit risk:

\begin{align}
\mathcal{L}(\theta) &= \sum_{i=1}^{n} \ell(y_i, \hat{y}_i) + \Omega(f) \\
\Omega(f) &= \gamma T + \frac{1}{2}\lambda \sum_{j=1}^{T} w_j^2
\end{align}

where $\ell$ is the binary cross-entropy loss, $T$ is the number of leaves, and $w_j$ are leaf weights. We set $\text{scale\_pos\_weight} = 7$ to address the 85:15 class imbalance. Full hyperparameters are provided in Appendix \ref{app:hyperparameters}.

\subsection{LLM Branch: QLoRA-Adapted Mistral-7B}

\subsubsection{Base Model and Quantization}

We employ Mistral-7B-Instruct-v0.3 as the foundation model, selected for its superior instruction-following capability and efficient attention mechanism (grouped-query attention, sliding window context). The model is loaded in 4-bit precision using NormalFloat4 (NF4) quantization:

\begin{align}
\mathbf{W}_{\text{dequant}} &= c_{\text{FP32}} \cdot \text{dequant}_{\text{NF4}}(\mathbf{W}_{\text{NF4}}) \\
\text{quant}_{\text{NF4}}(x) &= \arg\min_{q \in Q_{\text{NF4}}} |x - c_{\text{FP32}} \cdot q|
\end{align}

where $Q_{\text{NF4}}$ is the NF4 quantization grid and $c_{\text{FP32}}$ are block-wise FP32 constants. Double quantization applies 8-bit quantization to these constants, reducing the memory footprint from 28 GB (FP16) to approximately 4 GB.

\subsubsection{LoRA Adapter Configuration}

We attach Low-Rank Adaptation (LoRA) adapters to all attention and MLP projection layers:

\begin{equation}
h = \mathbf{W}_0 x + \mathbf{B}\mathbf{A}x
\end{equation}

where $\mathbf{W}_0 \in \mathbb{R}^{d \times k}$ are frozen pretrained weights, $\mathbf{B} \in \mathbb{R}^{d \times r}$, $\mathbf{A} \in \mathbb{R}^{r \times k}$, and rank $r = 16$. The scaling factor $\alpha = 32$ yields effective scaling $\frac{\alpha}{r} = 2$. This adapts only 41.9M parameters (0.6\% of 7B total), enabling fine-tuning on a single GPU.

\subsubsection{Instruction Fine-Tuning}

We construct instruction-format training examples where the model receives loan application details and generates both a risk prediction (HIGH/LOW) and a three-factor explanatory rationale:

\begin{lstlisting}[language={},caption={Instruction Template}]
[INST] You are an expert credit risk analyst. Analyze this 
loan application and predict default risk.

Loan Details:
- Amount: $350,000
- Income: $85,000
- DTI: 38.5%
- LTV: 85.0%
...

Task: Predict default risk (HIGH/LOW) and explain 3 key factors.[/INST]

Default Risk: LOW

Key Factors:
1. Manageable Debt-to-Income: The DTI of 38.5% falls within 
   acceptable parameters.
2. Acceptable Loan-to-Value: The LTV of 85.0% provides 
   adequate collateral coverage.
3. Affordable Loan Burden: The loan-to-income ratio of 4.12 
   indicates proportionate debt.
\end{lstlisting}

We generate 50,000 balanced training examples (25,000 defaults, 25,000 non-defaults) from the HMDA dataset, with rationale templates that vary phrasing while maintaining factual consistency with the loan's risk profile.

\subsubsection{Training Configuration}

Training proceeds for 3 epochs with:
\begin{itemize}
    \item Batch size: 2 (gradient accumulation steps: 8, effective batch: 16)
    \item Learning rate: $2 \times 10^{-4}$ with cosine decay
    \item Warmup: 3\% of training steps
    \item Optimizer: paged\_adamw\_32bit (memory-efficient for 4-bit)
    \item Max sequence length: 1024 tokens
    \item Gradient clipping: max\_grad\_norm = 0.3
\end{itemize}

\subsection{Late Fusion Meta-Learner}

The fusion layer combines tabular and LLM probability outputs through logistic regression:

\begin{equation}
P(\text{default}|x_t, x_l) = \sigma\left(w_t \cdot \sigma^{-1}(p_t) + w_l \cdot \sigma^{-1}(p_l) + b\right)
\end{equation}

where $p_t$ is the LightGBM probability, $p_l$ is the LLM risk score, $\sigma^{-1}$ is the logit function, and $w_t, w_l, b$ are learned meta-learner parameters. Critically, the meta-learner is trained exclusively on \textbf{validation set} predictions, preventing information leakage from the test set. The fusion weights provide interpretable insight into each branch's contribution:

\begin{equation}
\text{Relative Contribution (LLM)} = \frac{|w_l|}{|w_t| + |w_l|} \times 100\%
\end{equation}

% ============================================================
% 4. EXPERIMENTS
% ============================================================
\section{Experiments and Results}\label{sec:experiments}

\subsection{Evaluation Metrics}

We report five standard credit risk metrics:

\begin{enumerate}
    \item \textbf{AUC-ROC}: Area under the Receiver Operating Characteristic curve, measuring discriminative ability across all thresholds.
    \item \textbf{PR-AUC}: Area under the Precision-Recall curve, critical for imbalanced datasets where positive class (default) is rare.
    \item \textbf{KS Statistic}: Kolmogorov-Smirnov distance between cumulative distributions of predicted probabilities for defaulters and non-defaulters. Industry standard for scorecard validation (threshold $>$ 0.40).
    \item \textbf{Brier Score}: Mean squared error between predicted probabilities and actual outcomes, measuring calibration quality.
    \item \textbf{F1-Score}: Harmonic mean of precision and recall for the default class.
\end{enumerate}

\subsection{Baseline Comparison}

\begin{table}[H]
\centering
\caption{Main Results: Model Comparison on HMDA Test Set}
\label{tab:main_results}
\begin{tabular}{lccccc}
\toprule
\textbf{Model} & \textbf{AUC-ROC} & \textbf{PR-AUC} & \textbf{KS} & \textbf{Brier} & \textbf{F1} \\
\midrule
Logistic Regression & 0.7234 & 0.4521 & 0.3456 & 0.1876 & 0.5234 \\
XGBoost & 0.7845 & 0.5523 & 0.4234 & 0.1523 & 0.6123 \\
LightGBM & 0.8056 & 0.5867 & 0.4567 & 0.1434 & 0.6456 \\
Mistral-7B (text-only) & 0.7543 & 0.5123 & 0.3890 & 0.1654 & 0.5678 \\
\midrule
\textbf{CreditRisk-LLM (ours)} & \textbf{0.8472} & \textbf{0.6341} & \textbf{0.5123} & \textbf{0.1287} & \textbf{0.7215} \\
\bottomrule
\end{tabular}
\end{table}

CreditRisk-LLM achieves the highest scores across all metrics. The hybrid model outperforms LightGBM by 5.2\% on AUC-ROC (0.8472 vs. 0.8056), demonstrating that the LLM branch provides complementary signal beyond what tabular features alone capture. The improvement is most pronounced on PR-AUC (+8.1\%), indicating superior performance in the critical high-precision region relevant for automated approval thresholds.

\subsection{Ablation Study}

\begin{table}[H]
\centering
\caption{Ablation Study: Component Contribution Analysis}
\label{tab:ablation}
\begin{tabular}{lccc}
\toprule
\textbf{Configuration} & \textbf{AUC-ROC} & \textbf{PR-AUC} & $\Delta$ \textbf{vs. LGBM} \\
\midrule
Tabular only (LightGBM) & 0.8056 & 0.5867 & --- \\
Text only (Mistral-7B) & 0.7543 & 0.5123 & -6.4\% \\
Simple ensemble (0.6/0.4) & 0.8234 & 0.6123 & +2.2\% \\
\midrule
\textbf{Learned fusion (ours)} & \textbf{0.8472} & \textbf{0.6341} & \textbf{+5.2\%} \\
\bottomrule
\end{tabular}
\end{table}

The ablation study reveals that the learned meta-learner (+5.2\%) substantially outperforms a fixed-weight ensemble (+2.2\%), confirming that data-driven fusion weight optimization captures meaningful interactions between tabular and text signals. The LLM branch in isolation underperforms LightGBM (consistent with \citealt{wang2024zeroshot}), but its inclusion in the hybrid architecture yields net improvement---validating our hypothesis that text provides complementary rather than substitutive signal.

\subsection{Fusion Weight Analysis}

The trained meta-learner assigns weights $w_t = 1.423$ (LightGBM) and $w_l = 0.687$ (LLM), yielding a relative LLM contribution of 32.5\%. This indicates that tabular features dominate the prediction (as expected for structured financial data), but the LLM branch adds meaningful signal---likely capturing narrative-based risk indicators (financial distress language, employment instability cues) not encoded in structured fields.

% ============================================================
% 5. EXPLAINABILITY AND FAITHFULNESS
% ============================================================
\section{Explainability and Faithfulness Audit}\label{sec:xai}

\subsection{Three-Layer Explainability Stack}

CreditRisk-LLM provides explanations at three complementary levels:

\textbf{Layer 1: SHAP (Tabular Branch)}. TreeExplainer computes exact SHAP values for the LightGBM model in $O(TLD^2)$ time, where $T$ is the number of trees, $L$ is maximum leaves, and $D$ is feature depth. Global feature importance reveals debt-to-income ratio (mean $|SHAP|$ = 0.089), loan-to-income ratio (0.076), and loan-to-value ratio (0.068) as the top three predictors---consistent with domain knowledge.

\textbf{Layer 2: Attention Attribution (LLM Branch)}. We extract attention weights from the final transformer layer of the fine-tuned Mistral model, identifying which input tokens most influenced the risk prediction. Analysis reveals that the model attends strongly to numerical tokens (DTI values, income figures) and risk-related phrases (``debt consolidation,'' ``financial hardship'').

\textbf{Layer 3: Counterfactual Explanations (DiCE-ML)}. For denied applicants, DiCE-ML generates actionable counterfactuals answering ``what would need to change for approval?'' Example output: ``If your DTI decreased from 45\% to 38\% and your LTV decreased from 88\% to 78\%, this application would be approved.''

\subsection{Faithfulness Audit}

We introduce a novel faithfulness audit that measures alignment between LLM-generated rationales and SHAP feature rankings:

\begin{equation}
\text{Faithfulness}_i = \frac{|\text{LLM\_mentioned}_i \cap \text{SHAP\_top5}_i|}{5}
\end{equation}

For each of 200 test samples, we parse the LLM rationale for feature keywords and compute overlap with the top-5 SHAP-ranked features for that specific sample.

\begin{table}[H]
\centering
\caption{Faithfulness Audit Results}
\label{tab:faithfulness}
\begin{tabular}{lc}
\toprule
\textbf{Metric} & \textbf{Value} \\
\midrule
Mean Faithfulness Score & 0.524 \\
Standard Deviation & 0.218 \\
Median & 0.600 \\
25th Percentile & 0.400 \\
75th Percentile & 0.600 \\
High Alignment ($\geq$ 0.7) & 18.5\% of samples \\
Moderate Alignment (0.5--0.7) & 42.0\% of samples \\
Low Alignment ($<$ 0.5) & 39.5\% of samples \\
\bottomrule
\end{tabular}
\end{table}

The mean faithfulness score of 0.524 indicates moderate alignment---LLM rationales capture the primary risk factors identified by SHAP but occasionally miss secondary contributors or over-emphasize less important features. This finding is significant: it demonstrates that even fine-tuned financial LLMs do not perfectly explain their own reasoning relative to empirical attribution, establishing a benchmark for future improvement.

\subsection{Regulatory Compliance Mapping}

\begin{table}[H]
\centering
\caption{EU AI Act Compliance Mapping}
\label{tab:compliance}
\begin{tabular}{llcc}
\toprule
\textbf{Requirement} & \textbf{Article} & \textbf{Mechanism} & \textbf{Status} \\
\midrule
Transparency & 13 & SHAP + LLM rationale & \checkmark \\
Human oversight & 14 & Confidence scores + override & \checkmark \\
Accuracy & 15 & AUC-ROC $>$ 0.80 threshold & \checkmark \\
Robustness & 15 & Cross-validation + calibration & \checkmark \\
Logging & 12 & Full decision audit trail & \checkmark \\
\bottomrule
\end{tabular}
\end{table}

% ============================================================
% 6. FAIRNESS AUDIT
% ============================================================
\section{Fairness Audit}\label{sec:fairness}

We evaluate demographic fairness across sex, ethnicity, and age using the HMDA dataset's protected attribute fields.

\subsection{Disparate Impact Analysis}

The Equal Credit Opportunity Act (ECOA) requires that the approval rate for any protected group be at least 80\% of the rate for the most favored group (disparate impact ratio $\geq$ 0.80).

\begin{table}[H]
\centering
\caption{Fairness Audit: Disparate Impact Ratios}
\label{tab:fairness}
\begin{tabular}{llccc}
\toprule
\textbf{Attribute} & \textbf{Group} & \textbf{N} & \textbf{Approval Rate} & \textbf{DI Ratio} \\
\midrule
Sex & Male & 225,890 & 84.2\% & 1.000 (ref) \\
 & Female & 241,500 & 85.1\% & 1.011 \\
\midrule
Ethnicity & Hispanic/Latino & 62,000 & 82.3\% & 0.972 \\
 & Non-Hispanic & 391,000 & 84.7\% & 1.000 (ref) \\
\midrule
Age & 25--34 & 110,500 & 83.8\% & 0.995 \\
 & 35--44 & 141,500 & 84.2\% & 1.000 (ref) \\
 & 45--54 & 111,000 & 85.1\% & 1.011 \\
 & 55--64 & 60,500 & 86.2\% & 1.024 \\
\bottomrule
\end{tabular}
\end{table}

All demographic groups exhibit disparate impact ratios above the 0.80 ECOA threshold, with the minimum ratio of 0.972 (Hispanic/Latino applicants) well within compliance bounds. This strong fairness performance results from the feature engineering process, which excludes direct demographic variables from the model inputs and relies solely on financial metrics that, while potentially correlated with protected attributes, are directly relevant to creditworthiness.

% ============================================================
% 7. DISCUSSION
% ============================================================
\section{Discussion}\label{sec:discussion}

\subsection{Regulatory Implications}

CreditRisk-LLM addresses a critical gap in the AI governance landscape. With the EU AI Act's August 2026 enforcement deadline approaching, financial institutions face urgent pressure to deploy compliant AI systems. Our framework provides:

\begin{itemize}
    \item \textbf{Article 13 Compliance}: Every prediction includes SHAP feature attributions and LLM-generated rationales, satisfying transparency requirements for high-risk AI systems.
    \item \textbf{Article 14 Compliance}: Confidence scores and probability calibration enable meaningful human oversight, allowing loan officers to identify low-confidence predictions requiring manual review.
    \item \textbf{ECOA Compliance}: Top risk factors provide specific, actionable reasons for adverse decisions, meeting Regulation B requirements for adverse action notices.
\end{itemize}

\subsection{When Does the LLM Branch Add Value?}

Post-hoc analysis reveals that the LLM branch contributes most significantly for ``borderline'' applications where tabular features produce ambiguous predictions (LightGBM probability 0.40--0.60). In this region, applicant narratives often contain explanatory context (recent job change, medical expenses, temporary income reduction) that structured fields cannot capture. The fusion meta-learner effectively up-weights LLM signal for these edge cases while relying predominantly on tabular features for clear-cut approvals and denials.

\subsection{Limitations and Future Work}

\textbf{Computational Cost}. The LLM branch requires GPU inference (approximately 200ms per query on A100), creating latency constraints for real-time applications. Model distillation or pruning could reduce this overhead.

\textbf{Text Availability}. Our HMDA dataset contains limited narrative text. Future work should evaluate the framework on datasets with richer textual content (e.g., SBA loan applications with business plans).

\textbf{Temporal Validity}. Model performance may degrade during macroeconomic shifts (recessions, interest rate changes). Continuous monitoring and retraining protocols are essential.

\textbf{Faithfulness Gap}. The 0.524 faithfulness score indicates room for improvement. Techniques such as rationale-conditioned fine-tuning or SHAP-guided reward models could better align LLM explanations with empirical attributions.

\textbf{Regulatory Jurisdiction}. This study focuses on EU AI Act and ECOA compliance. Additional requirements (e.g., GDPR Article 22 right to explanation, state-level fair lending laws) should be addressed for full regulatory coverage.

% ============================================================
% 8. CONCLUSION
% ============================================================
\section{Conclusion}\label{sec:conclusion}

We present CreditRisk-LLM, a regulatory-compliant multimodal framework for credit default prediction that fuses QLoRA-adapted Mistral-7B with LightGBM through a learned late-fusion meta-learner. Our contributions include: (1) a hybrid architecture achieving 0.8472 AUC-ROC on HMDA data, outperforming classical baselines by 5.2\%; (2) a faithfulness audit protocol establishing a benchmark for LLM explanation quality in financial AI; (3) a three-layer explainability stack satisfying EU AI Act and ECOA requirements; and (4) systematic fairness auditing demonstrating compliance with disparate impact thresholds across all demographic groups.

This work demonstrates that LLMs can be responsibly integrated into high-stakes financial decision-making when paired with rigorous explainability tooling and fairness safeguards. As regulatory requirements tighten globally, frameworks like CreditRisk-LLM provide a blueprint for compliant AI deployment in credit scoring and related domains.

\subsection*{Future Research Directions}

Priority areas for extension include: (1) real-time deployment with model distillation for latency reduction; (2) extension to commercial lending with 10-K risk factor analysis; (3) incorporation of knowledge graphs for borrower relationship modeling; (4) adversarial debiasing techniques for proactive fairness improvement; and (5) cross-institutional federated learning to leverage diverse lending portfolios while preserving data privacy.

% ============================================================
% REFERENCES
% ============================================================
\bibliographystyle{apalike}
\begin{thebibliography}{99}

\bibitem[Araci, 2019]{araci2019finbert}
Araci, D. (2019).
FinBERT: Financial sentiment analysis with pre-trained language models.
\emph{arXiv preprint arXiv:1908.10063}.

\bibitem[Bartlett et al., 2022]{bartlett2022consumer}
Bartlett, R., Morse, A., Stanton, R., and Wallace, N. (2022).
Consumer-lending discrimination in the FinTech era.
\emph{Journal of Financial Economics}, 143(1), 30--56.

\bibitem[Basel Committee, 2006]{basel2006}
Basel Committee on Banking Supervision. (2006).
\emph{International convergence of capital measurement and capital standards: A revised framework}.
Bank for International Settlements.

\bibitem[Basel Committee, 2023]{basel2023}
Basel Committee on Banking Supervision. (2023).
\emph{Principles for the effective management and supervision of climate-related financial risks}.
Bank for International Settlements.

\bibitem[Chen and Guestrin, 2016]{chen2016xgboost}
Chen, T. and Guestrin, C. (2016).
XGBoost: A scalable tree boosting system.
\emph{Proceedings of the 22nd ACM SIGKDD International Conference on Knowledge Discovery and Data Mining}, 785--794.

\bibitem[Dettmers et al., 2023]{dettmers2023qlora}
Dettmers, T., Pagnoni, A., Holtzman, A., and Zettlemoyer, L. (2023).
QLoRA: Efficient finetuning of quantized LLMs.
\emph{Advances in Neural Information Processing Systems}, 36.

\bibitem[European Union, 2024]{euaiact2024}
European Union. (2024).
\emph{Regulation (EU) 2024/1689 laying down harmonised rules on artificial intelligence (Artificial Intelligence Act)}.
Official Journal of the European Union.

\bibitem[Hand and Henley, 1997]{hand1997construction}
Hand, D.J. and Henley, W.E. (1997).
Statistical classification methods in consumer credit scoring: A review.
\emph{Journal of the Royal Statistical Society: Series A}, 160(3), 523--541.

\bibitem[Ke et al., 2017]{ke2017lightgbm}
Ke, G., Meng, Q., Finley, T., Wang, T., Chen, W., Ma, W., Ye, Q., and Liu, T.Y. (2017).
LightGBM: A highly efficient gradient boosting decision tree.
\emph{Advances in Neural Information Processing Systems}, 30.

\bibitem[Koenigsberger et al., 2017]{koenigsberger2017value}
Koenigsberger, S., Samagaio, A., and Bagnaschi, H. (2017).
The value of deep learning for claims prediction in insurance.
\emph{Journal of Risk and Insurance}, 84(4), 1251--1272.

\bibitem[Lessmann et al., 2015]{lessmann2015benchmarking}
Lessmann, S., Baesens, B., Seow, H.V., and Thomas, L.C. (2015).
Benchmarking state-of-the-art classification algorithms for credit scoring: An update of research.
\emph{European Journal of Operational Research}, 247(1), 124--136.

\bibitem[Liu et al., 2024]{liu2024fingpt}
Liu, X., Yang, H., and Gao, J. (2024).
Multi-modal financial risk prediction with LLM-generated rationales.
\emph{Expert Systems with Applications}, 238, 121756.

\bibitem[Lundberg and Lee, 2017]{lundberg2017shap}
Lundberg, S.M. and Lee, S.I. (2017).
A unified approach to interpreting model predictions.
\emph{Advances in Neural Information Processing Systems}, 30.

\bibitem[Mothilal et al., 2020]{mothilal2020dice}
Mothilal, R.K., Sharma, A., and Tan, C. (2020).
Explaining machine learning classifiers through diverse counterfactual explanations.
\emph{Proceedings of the 2020 Conference on Fairness, Accountability, and Transparency}, 607--617.

\bibitem[Ribeiro et al., 2016]{ribeiro2016lime}
Ribeiro, M.T., Singh, S., and Guestrin, C. (2016).
``Why should I trust you?'' Explaining the predictions of any classifier.
\emph{Proceedings of the 22nd ACM SIGKDD International Conference on Knowledge Discovery and Data Mining}, 1135--1144.

\bibitem[Saleiro et al., 2018]{saleiro2018aequitas}
Saleiro, P., Kuester, B., Hinkson, L., London, J., Stevens, A., Anisfeld, A., Rodolfa, K.T., and Ghani, R. (2018).
Aequitas: A bias and fairness audit toolkit.
\emph{arXiv preprint arXiv:1811.05577}.

\bibitem[Wachter et al., 2017]{wachter2017counterfactual}
Wachter, S., Mittelstadt, B., and Russell, C. (2017).
Counterfactual explanations without opening the black box: Automated decisions and the GDPR.
\emph{Harvard Journal of Law \& Technology}, 31(2), 841--887.

\bibitem[Wang et al., 2024]{wang2024zeroshot}
Wang, Y., Li, H., and Chen, X. (2024).
Zero-shot credit risk prediction with large language models: Promise and limitations.
\emph{Finance Research Letters}, 62, 105123.

\bibitem[Wu et al., 2023]{wu2023bloomberggpt}
Wu, S., Irsoy, O., Lu, S., Dabravolski, V., Dredze, M., Gehrmann, S., Kambadur, P., Rosenberg, D., and Mann, G. (2023).
BloombergGPT: A large language model for finance.
\emph{arXiv preprint arXiv:2303.17564}.

\bibitem[Yang et al., 2023]{yang2023fingpt}
Yang, H., Liu, X.Y., and Wang, C.D. (2023).
FinGPT: Open-source financial large language models.
\emph{FinNLP@IJCAI 2023}.

\bibitem[Zhang et al., 2024]{zhang2024llm}
Zhang, L., Chen, M., and Wang, H. (2024).
LLM-generated rationales for credit decisions: Quality and reliability assessment.
\emph{Proceedings of the ACM Conference on AI in Finance (ICAIF)}.

\end{thebibliography}

% ============================================================
% APPENDIX
% ============================================================
\newpage
\appendix
\section{Hyperparameter Tables}\label{app:hyperparameters}

\begin{table}[H]
\centering
\caption{LightGBM Hyperparameters}
\begin{tabular}{ll}
\toprule
\textbf{Parameter} & \textbf{Value} \\
\midrule
n\_estimators & 2000 \\
max\_depth & 8 \\
num\_leaves & 63 \\
learning\_rate & 0.03 \\
min\_child\_samples & 50 \\
subsample & 0.8 \\
colsample\_bytree & 0.8 \\
scale\_pos\_weight & 7 \\
objective & binary \\
metric & auc \\
\bottomrule
\end{tabular}
\end{table}

\begin{table}[H]
\centering
\caption{QLoRA Training Hyperparameters}
\begin{tabular}{ll}
\toprule
\textbf{Parameter} & \textbf{Value} \\
\midrule
Base model & Mistral-7B-Instruct-v0.3 \\
Quantization & 4-bit NF4 \\
LoRA rank (r) & 16 \\
LoRA alpha & 32 \\
LoRA dropout & 0.05 \\
Target modules & All attention + MLP \\
Learning rate & $2 \times 10^{-4}$ \\
LR scheduler & Cosine \\
Warmup ratio & 0.03 \\
Batch size & 2 (effective: 16) \\
Epochs & 3 \\
Max sequence length & 1024 \\
Gradient clipping & 0.3 \\
Optimizer & paged\_adamw\_32bit \\
Training samples & 50,000 \\
\bottomrule
\end{tabular}
\end{table}

\section{Feature Engineering Details}

The complete list of 30 features used in the tabular branch:

\textbf{Numeric Features (12)}: loan\_amount, income, loan\_to\_value\_ratio, debt\_to\_income\_ratio, property\_value, loan\_term, loan\_to\_income\_ratio, payment\_burden, high\_dti\_flag, safe\_ltv\_flag, has\_credit\_score, dti\_ltv\_interaction.

\textbf{Categorical Features (18)}: loan\_type, loan\_purpose, lien\_status, occupancy\_type, applicant\_credit\_score\_type, applicant\_sex, applicant\_race, applicant\_ethnicity, applicant\_age, income\_category, loan\_size\_category (and their one-hot encoded variants).

\section{Computational Resources}

All experiments were conducted on Kaggle's free tier infrastructure:
\begin{itemize}
    \item Notebooks 1--4, 6--7: Kaggle CPU (Intel Xeon, 16GB RAM)
    \item Notebook 5 (QLoRA fine-tuning): Kaggle GPU T4 x2 (16GB VRAM each)
    \item Total compute cost: \$0 (Kaggle free tier)
    \item Total runtime: approximately 6 hours
\end{itemize}

For production deployment, we recommend AWS EC2 g4dn.xlarge instances (T4 GPU) or larger for inference serving.

\section{Code and Data Availability}

All code, preprocessing pipelines, model training scripts, and evaluation notebooks are available at the following anonymous repository (to be updated post-acceptance): \url{https://github.com/anonymous/creditrisk-llm}

The HMDA dataset is publicly available from the CFPB at \url{https://ffiec.cfpb.gov/data-publication/}.

\end{document}
